\chapter{Charakterisierungen und Randfälle}
\label{b07guide:ch:criteria}

\section{Volles Bild und nichtleere Fasern}

\begin{theorem}[Bild- und Faserkriterien]
\label{b07guide:thm:pointcriteria}
Für eine Funktion $f\colon A\to B$ sind äquivalent:
\begin{enumerate}
\item $f$ ist surjektiv auf $B$.
\item $f[A]=B$.
\item Für jedes $b\in B$ gilt $\operatorname{Fib}_f(b)\ne\varnothing$.
\item Für jedes nichtleere $V\subseteq B$ gilt
      $f^{-1}[V]\ne\varnothing$.
\item Es gibt keine echte Teilmenge $D\subsetneq B$,
      für die derselbe Graph eine Funktion $f\colon A\to D$ ist.
\end{enumerate}
In Bedingung~4 genügt es, alle Einermengen $V=\{b\}$ zu prüfen.
\end{theorem}

\begin{theorem}[Kriterium für Nichtsurjektivität]
\label{b07guide:thm:failure}
Für $f\colon A\to B$ sind äquivalent:
\[
 \begin{aligned}
 f\text{ ist nicht surjektiv auf }B
 &\Longleftrightarrow B\setminus f[A]\ne\varnothing\\
 &\Longleftrightarrow
   \exists b\in B\;\forall a\in A\quad f(a)\ne b.
 \end{aligned}
\]
\end{theorem}

\section{Charakterisierung durch Teilmengen}

\begin{definition}[Bildfunktion auf der Potenzmenge]
Für $f\colon A\to B$ sei
\[
 f_*\colon\mathcal P(A)\to\mathcal P(B),
 \qquad U\mapsto f[U].
\]
\end{definition}

\begin{theorem}[Rückgewinnung sämtlicher Zielteilmengen]
\label{b07guide:thm:setcriteria}
Für $f\colon A\to B$ sind äquivalent:
\begin{enumerate}
\item $f$ ist surjektiv auf $B$.
\item Für jedes $V\subseteq B$ gilt $f[f^{-1}[V]]=V$.
\item Jede Teilmenge von $B$ ist das Bild einer Teilmenge von $A$;
      also ist $f_*$ surjektiv.
\item Für alle $V,W\subseteq B$ gilt
      $f^{-1}[V]=f^{-1}[W]\Rightarrow V=W$.
\item Für alle $V,W\subseteq B$ gilt
      $f^{-1}[V]\subseteq f^{-1}[W]\Rightarrow V\subseteq W$.
\end{enumerate}
In Bedingung~2 genügt bereits $V=B$.
\end{theorem}

\section{Leere und kleine Zielmengen}

\begin{theorem}[Leere Bereiche]
\label{b07guide:thm:empty}
Es gelten
\[
 \operatorname{Sur}(\varnothing,B)=
 \begin{cases}
 \{\varnothing\},& B=\varnothing,\\
 \varnothing,& B\ne\varnothing,
 \end{cases}
 \qquad
 \operatorname{Sur}(A,\varnothing)=
 \begin{cases}
 \{\varnothing\},& A=\varnothing,\\
 \varnothing,& A\ne\varnothing.
 \end{cases}
\]
Insbesondere ist die leere Funktion
$\varnothing\colon\varnothing\twoheadrightarrow\varnothing$
surjektiv. Für jede Surjektion $f\colon A\twoheadrightarrow B$ gilt
\[
 A=\varnothing\quad\Longleftrightarrow\quad B=\varnothing.
\]
\end{theorem}

\begin{theorem}[Einermengen und konstante Funktionen]
\label{b07guide:thm:constant}
Für jede Menge $A$ und jede Einermenge $\{b\}$ gilt
\[
 \operatorname{Sur}(A,\{b\})\ne\varnothing
 \quad\Longleftrightarrow\quad A\ne\varnothing.
\]
Bei $A\ne\varnothing$ ist die konstante Funktion $a\mapsto b$
die einzige Funktion und damit die einzige Surjektion nach $\{b\}$.

Allgemeiner sei $b\in B$ und $c\colon A\to B$ die konstante
Funktion mit Wert $b$. Dann gilt
\[
 c\text{ surjektiv auf }B
 \quad\Longleftrightarrow\quad
 A\ne\varnothing\ \text{ und }\ B=\{b\}.
\]
Eine Funktion mit einer Einermenge als Definitionsbereich ist
genau dann surjektiv, wenn ihre Zielmenge ebenfalls eine Einermenge ist.
\end{theorem}